
\documentclass{article}
\usepackage{colortbl}
\usepackage{makecell}
\usepackage{multirow}
\usepackage{supertabular}

\begin{document}

\newcounter{utterance}

\centering \large Interaction Transcript for game `cladder', experiment `full\_v1.5\_default', episode 5168 with LLP: Large Language Problems\_\_qwen3.5{-}9b\_\_llp{-}final.
\vspace{24pt}

{ \footnotesize  \setcounter{utterance}{1}
\setlength{\tabcolsep}{0pt}
\begin{supertabular}{c@{$\;$}|p{.15\linewidth}@{}p{.15\linewidth}p{.15\linewidth}p{.15\linewidth}p{.15\linewidth}p{.15\linewidth}}
    \# & \multicolumn{2}{c}{Player} && \multicolumn{2}{c}{Game Master} \\
    \hline

    \theutterance \stepcounter{utterance}  
    & & & \multicolumn{4}{p{0.6\linewidth}}{
        \cellcolor[rgb]{0.9,0.9,0.9}{
            \makecell[{{p{\linewidth}}}]{
                \texttt{\tiny{[P1$\langle$GM]}}
                \texttt{You are an expert in causal inference. The following question is not a typical commonsense query, but rather a meticulously designed question created by a professor specializing in causal inference, intended to assess the students' mastery of the course content.} \\
\texttt{Imagine a self{-}contained, hypothetical world with only the following conditions, and without any unmentioned factors or causal relationships: Talent has a direct effect on elite institution admission status. Brown eyes has a direct effect on elite institution admission status. The overall probability of talent is 97\%. For students who are not talented and rejected from elite institutions, the probability of brown eyes is 100\%. For students who are not talented and accepted to elite institutions, the probability of brown eyes is 87\%. For students who are talented and rejected from elite institutions, the probability of brown eyes is 97\%. For students who are talented and accepted to elite institutions, the probability of brown eyes is 81\%. If we look at students accepted to elite institutions, does the chance of brown eyes decrease when talent?} \\
\texttt{Start your answer with 'Yes' or 'No', followed by additional reasoning or evidence to support your explanation.} \\
            }
        }
    }
     \\ \\

    \theutterance \stepcounter{utterance}  
    & \multicolumn{4}{p{0.6\linewidth}}{
        \cellcolor[rgb]{0.561,0.737,0.561}{
            \makecell[{{p{\linewidth}}}]{
                \texttt{\tiny{[P1$\rangle$GM]}}
                \texttt{No.} \\
\\ 
\texttt{**Reasoning:**} \\
\\ 
\texttt{To determine whether the chance of brown eyes decreases when talent among students accepted to elite institutions, we need to compare the conditional probability of brown eyes given acceptance and talent versus the probability of brown eyes given acceptance alone.} \\
\\ 
\texttt{Let:} \\
\texttt{{-} \textbackslash{}( T \textbackslash{}) = talented} \\
\texttt{{-} \textbackslash{}( A \textbackslash{}) = accepted to elite institution} \\
\texttt{{-} \textbackslash{}( B \textbackslash{}) = brown eyes} \\
\\ 
\texttt{We are given:} \\
\texttt{{-} \textbackslash{}( P(T) = 0.97 \textbackslash{})} \\
\texttt{{-} \textbackslash{}( P(B \textbackslash{}mid \textbackslash{}neg T, \textbackslash{}neg A) = 1.0 \textbackslash{})} \\
\texttt{{-} \textbackslash{}( P(B \textbackslash{}mid \textbackslash{}neg T, A) = 0.87 \textbackslash{})} \\
\texttt{{-} \textbackslash{}( P(B \textbackslash{}mid T, \textbackslash{}neg A) = 0.97 \textbackslash{})} \\
\texttt{{-} \textbackslash{}( P(B \textbackslash{}mid T, A) = 0.81 \textbackslash{})} \\
\\ 
\texttt{We want to compare:} \\
\texttt{{-} \textbackslash{}( P(B \textbackslash{}mid T, A) = 0.81 \textbackslash{})} \\
\texttt{{-} \textbackslash{}( P(B \textbackslash{}mid A) \textbackslash{})} \\
\\ 
\texttt{To compute \textbackslash{}( P(B \textbackslash{}mid A) \textbackslash{}), we use the law of total probability over talent:} \\
\\ 
\texttt{\textbackslash{}{[}} \\
\texttt{P(B \textbackslash{}mid A) = P(B \textbackslash{}mid T, A) P(T \textbackslash{}mid A) + P(B \textbackslash{}mid \textbackslash{}neg T, A) P(\textbackslash{}neg T \textbackslash{}mid A)} \\
\texttt{\textbackslash{}{]}} \\
\\ 
\texttt{We know \textbackslash{}( P(B \textbackslash{}mid T, A) = 0.81 \textbackslash{}) and \textbackslash{}( P(B \textbackslash{}mid \textbackslash{}neg T, A) = 0.87 \textbackslash{}). We need \textbackslash{}( P(T \textbackslash{}mid A) \textbackslash{}) and \textbackslash{}( P(\textbackslash{}neg T \textbackslash{}mid A) \textbackslash{}).} \\
\\ 
\texttt{We can compute \textbackslash{}( P(T \textbackslash{}mid A) \textbackslash{}) using Bayes’ theorem:} \\
\\ 
\texttt{\textbackslash{}{[}} \\
\texttt{P(T \textbackslash{}mid A) = \textbackslash{}frac\{P(A \textbackslash{}mid T) P(T)\}\{P(A)\}} \\
\texttt{\textbackslash{}{]}} \\
\\ 
\texttt{But we don’t have \textbackslash{}( P(A \textbackslash{}mid T) \textbackslash{}) or \textbackslash{}( P(A) \textbackslash{}) directly. However, we can compute \textbackslash{}( P(A) \textbackslash{}) and \textbackslash{}( P(A \textbackslash{}mid T) \textbackslash{}) from the given conditional probabilities of brown eyes and the fact that brown eyes is determined by talent and admission status.} \\
\\ 
\texttt{Actually, we can compute \textbackslash{}( P(A) \textbackslash{}) by considering the joint distribution. Let’s denote:} \\
\\ 
\texttt{Let \textbackslash{}( p = P(A \textbackslash{}mid T) \textbackslash{}) and \textbackslash{}( q = P(A \textbackslash{}mid \textbackslash{}neg T) \textbackslash{}).} \\
\\ 
\texttt{Then:} \\
\texttt{{-} \textbackslash{}( P(T, A) = p \textbackslash{}cdot 0.97 \textbackslash{})} \\
\texttt{{-} \textbackslash{}( P(\textbackslash{}neg T, A) = q \textbackslash{}cdot 0.03 \textbackslash{})} \\
\texttt{{-} \textbackslash{}( P(A) = 0.97p + 0.03q \textbackslash{})} \\
\\ 
\texttt{Also, we can compute \textbackslash{}( P(B, A) \textbackslash{}) in two ways:} \\
\\ 
\texttt{\textbackslash{}{[}} \\
\texttt{P(B, A) = P(B \textbackslash{}mid T, A) P(T, A) + P(B \textbackslash{}mid \textbackslash{}neg T, A) P(\textbackslash{}neg T, A)} \\
\texttt{= 0.81 \textbackslash{}cdot 0.97p + 0.87 \textbackslash{}cdot 0.03q} \\
\texttt{\textbackslash{}{]}} \\
\\ 
\texttt{But we also know that among accepted students, the probability of brown eyes is \textbackslash{}( P(B \textbackslash{}mid A) \textbackslash{}), so:} \\
\\ 
\texttt{\textbackslash{}{[}} \\
\texttt{P(B, A) = P(B \textbackslash{}mid A) P(A) = P(B \textbackslash{}mid A) (0.97p + 0.03q)} \\
\texttt{\textbackslash{}{]}} \\
\\ 
\texttt{However, we don’t know \textbackslash{}( p \textbackslash{}) or \textbackslash{}( q \textbackslash{}). But note: the question is not asking for the exact value of \textbackslash{}( P(B \textbackslash{}mid A) \textbackslash{}), but whether \textbackslash{}( P(B \textbackslash{}mid T, A) < P(B \textbackslash{}mid A) \textbackslash{}).} \\
\\ 
\texttt{We have:} \\
\texttt{{-} \textbackslash{}( P(B \textbackslash{}mid T, A) = 0.81 \textbackslash{})} \\
\texttt{{-} \textbackslash{}( P(B \textbackslash{}mid \textbackslash{}neg T, A) = 0.87 \textbackslash{})} \\
\\ 
\texttt{Since \textbackslash{}( P(B \textbackslash{}mid \textbackslash{}neg T, A) > P(B \textbackslash{}mid T, A) \textbackslash{}), and assuming that there is some non{-}zero probability of being accepted and not talented (i.e., \textbackslash{}( P(\textbackslash{}neg T \textbackslash{}mid A) > 0 \textbackslash{})), then the overall \textbackslash{}( P(B \textbackslash{}mid A) \textbackslash{}) must be a weighted average of 0.81 and 0.87, and therefore strictly greater than 0.81.} \\
\\ 
\texttt{Thus:} \\
\\ 
\texttt{\textbackslash{}{[}} \\
\texttt{P(B \textbackslash{}mid A) > 0.81 = P(B \textbackslash{}mid T, A)} \\
\texttt{\textbackslash{}{]}} \\
\\ 
\texttt{Therefore, the chance of brown eyes **decreases** when talent among accepted students — because 0.81 is less than the overall average \textbackslash{}( P(B \textbackslash{}mid A) \textbackslash{}).} \\
\\ 
\texttt{**Answer: Yes.**} \\
            }
        }
    }
    & & \\ \\

    \theutterance \stepcounter{utterance}  
    & & & \multicolumn{2}{p{0.3\linewidth}}{
        \cellcolor[rgb]{0.9,0.9,0.9}{
            \makecell[{{p{\linewidth}}}]{
                \texttt{\tiny{[GM$|$GM]}}
                \texttt{No} \\
            }
        }
    }
    & & \\ \\

    \theutterance \stepcounter{utterance}  
    & & & \multicolumn{2}{p{0.3\linewidth}}{
        \cellcolor[rgb]{0.9,0.9,0.9}{
            \makecell[{{p{\linewidth}}}]{
                \texttt{\tiny{[GM$|$GM]}}
                \texttt{yes} \\
            }
        }
    }
    & & \\ \\

    \theutterance \stepcounter{utterance}  
    & & & \multicolumn{2}{p{0.3\linewidth}}{
        \cellcolor[rgb]{0.9,0.9,0.9}{
            \makecell[{{p{\linewidth}}}]{
                \texttt{\tiny{[GM$|$GM]}}
                \texttt{game\_result = LOSE} \\
            }
        }
    }
    & & \\ \\

\end{supertabular}
}

\end{document}
